\documentclass[12pt,a4paper]{article}
\usepackage[russian]{babel}
\usepackage{fontspec}
\usepackage{geometry}
\usepackage{booktabs}
\usepackage{longtable}
\usepackage{array}
\usepackage{multirow}
\usepackage{enumitem}
\usepackage{titlesec}
\usepackage{xcolor}
\usepackage{siunitx}

\usepackage{pdfpages}

\usepackage{tablefootnote}

% Дополнительные пакеты для данного документа
\usepackage{caption}
\usepackage{subcaption}
\usepackage{listings}
% \usepackage{hyperref}

\geometry{a4paper, margin=25mm}
\setmainfont{Liberation Serif}
\setsansfont{Liberation Sans}

\definecolor{adblue}{RGB}{0,84,166}

\titleformat{\section}{\large\bfseries\color{adblue}}{}{0em}{}[\titlerule]
\titleformat{\subsection}{\normalsize\bfseries}{}{0em}{}
\titleformat{\subsubsection}{\normalsize\itshape}{}{0em}{}

\sisetup{output-decimal-marker={,}}

\usepackage{tikz}
\usetikzlibrary{positioning, calc, backgrounds, math, 3d, perspective, arrows.meta, graphs, graphdrawing, fit, shapes.geometric, trees}
    

\usepackage{pgfplots} % для построения графиков
\pgfplotsset{compat=1.18} % версия для совместимости
\usepackage{tkz-euclide}

% --- Подключение пакета алгоритмов (без флага algo2e, чтобы вернуть \begin{algorithm}) ---
\usepackage[ruled,vlined]{algorithm2e}
%\usepackage{caption}
\newcommand{\Gray}[1]{\textcolor{gray}{#1}} 
% Настройка подписей
\DeclareCaptionLabelFormat{gost}{#1 #2}
\captionsetup[table]{labelformat=gost,labelsep=endash,justification=justified,singlelinecheck=false,font=normal}
\captionsetup[figure]{labelformat=gost,labelsep=endash,justification=centering,font=normal}

% Локализация algorithm2e на русский (ОБЯЗАТЕЛЬНО после загрузки пакета)
\SetAlgorithmName{Алгоритм}{}{}
\SetAlgoCaptionSeparator{---}
\SetKwInput{Input}{Вход}
\SetKwInput{Output}{Выход}

\usepackage{url}
% Говорим пакету использовать моноширинный шрифт (как у texttt)
\Urlmuskip=0mu plus 1mu % гибкие пробелы для формул, если нужно
\DeclareUrlCommand\code{\urlstyle{tt}} 

% Библиотека схемотехники
\usepackage[siunitx, RPvoltages, american]{circuitikzgit}
\ctikzset{
    amplifiers/fill=cyan,
    sources/fill=green!20,
    diodes/fill=white,
    resistors/fill=violet,
    grounds/scale=1.2,
}

\usepackage{iftex}

\ifXeTeX % Если используется TeTeX или TeLaTeX
  \typeout{Режим XeTeX}
  \usepackage{xltxtra}
  \usepackage{fontspec}
  \usepackage{polyglossia}
  \setmainlanguage{russian}
  \setotherlanguage{english}
  \setkeys{russian}{babelshorthands=true} % По идее включает <<, >>, -- ---

  \setmainfont{Times New Roman}[Ligatures=TeX]
  \setromanfont{Times New Roman}[Ligatures=TeX]
  \setsansfont{Arial}[Ligatures=TeX]
  \setmonofont{Courier New}[Ligatures=TeX] 

  \newfontfamily{\cyrillicfont}{Times New Roman}
  \newfontfamily{\cyrillicfontrm}{Times New Roman}
  \newfontfamily{\cyrillicfonttt}{Courier New}
  \newfontfamily{\cyrillicfontsf}{Arial}
\fi

\ifLuaTeX % Если используется LuaLatex
  \typeout{Режим LuaTeX}
%  \usepackage{luatextra}
  \usepackage{fontspec}
  \defaultfontfeatures{Ligatures={TeX}}
  \usepackage{polyglossia}
  \setmainlanguage{russian}
  \setotherlanguage{english}
  \setmainfont{Times New Roman}[Script=Cyrillic, Ligatures=TeX]
  \setromanfont{Times New Roman}[Script=Cyrillic, Ligatures=TeX]
  \setsansfont{Arial}[Script=Cyrillic, Ligatures=TeX]
  \setmonofont{Courier New}[Ligatures=TeX]
  \newfontfamily{\cyrillicfont}{Times New Roman}[Script=Cyrillic, Ligatures=TeX]
  \newfontfamily{\cyrillicfontrm}{Times New Roman}[Script=Cyrillic, Ligatures=TeX]
  \newfontfamily{\cyrillicfonttt}{Courier New}
  \newfontfamily{\cyrillicfontsf}{Arial}[Script=Cyrillic, Ligatures=TeX]

  \usegdlibrary{trees}
  \usegdlibrary{force}
\fi

\ifPDFTeX % Если используется 8битный PDFTex
  \typeout{8-битный режим pdfTex}
  % For pdfTeX we must use old
  % 8-bit font technologies
  \usepackage[T2A]{fontenc}
  \usepackage[english,russian]{babel}
  \usepackage[utf8]{inputenc}
  \usepackage[english, russian]{babel}

  %\usepackage{newtxmath}
  \usepackage{amsmath}
  \usepackage{amsthm}
  \usepackage{amssymb}
  \usepackage{amsfonts}
  \usepackage{mathtext}
\fi

\usepackage{amsthm}  % Возможности AMSMath
\usepackage{amsmath}
\usepackage{amssymb}

\usepackage{esvect} % Для красивых стрелок
\usepackage{bm} % Для жирных символов

%Hyphenation rules
\usepackage{hyphenat}
\hyphenation{ма-те-ма-ти-ка вос-ста-нав-ли-вать ото-бра-жа-ют-ся сге-не-ри-ро-ван-ном сим-во-лы}
% \usepackage[english, russian]{babel}

\usepackage{graphicx}
\graphicspath{ {./images/} }

\usepackage{empheq}

\usepackage{float}

% \usepackage{comment}

\usepackage[hidelinks,
    unicode=true,          % поддерживает Unicode в закладках
    psdextra=true,      % автоматически конвертирует простую математику в текст
    % focusdefault=false
]{hyperref}


% Окружение "Теорема"
\newtheorem{theorem}{Теорема}[section]
\newtheorem{corollary}{Corollary}[theorem]
\newtheorem{lemma}[theorem]{Lemma}

% Окружение "Определение"
\theoremstyle{definition} % Прямой шрифт, нумерация
\newtheorem{definition}{Определение}[section]
\newtheorem{example}{Пример}[section]

% Окружение "Замечание"
\theoremstyle{remark} % Прямой шрифт, без нумерации
% Окуржение "Решение"
\newtheorem*{solution}{Решение}
\newtheorem*{remark}{Замечание}

\usepackage{polynom}
\usepackage{tabularx}
\usepackage{cancel}

\addto\captionsrussian{%
  \renewcommand{\figurename}{Рис.}%
  \renewcommand{\tablename}{Табл.}%
}

\renewcommand{\arraystretch}{1.5} % Высота строки таблицы

\renewcommand{\labelitemi}{\(\circ\)} % Бюллет -- пустой кружок в itemize


\begin{document}

\section*{ТЕОРЕТИЧЕСКИЙ АНАЛИЗ ПРОСТРАНСТВЕННОЙ ДИСКРЕТИЗАЦИИ И НАУЧНАЯ НОВИЗНА МЕТОДА}

\subsection*{1. Физико-математическая модель формирования мозаичного сигнала сенсора}

Процесс регистрации оптического излучения твердотельным полупроводниковым сенсором фотокамеры Nikon D600 (шаблон маски Байера --- RGGB) может быть формализован как последовательное прохождение непрерывного двумерного светового распределения сквозь оптический тракт и пространственно-модулирующую дискретизирующую структуру.

Пусть $I_{\text{orig}}(x, y)$ --- идеальное непрерывное полутоновое изображение сцены. Прохождение светового потока через реальную оптическую систему (объектив) эквивалентно операции двумерной свёртки с функцией размытия точки (PSF --- \textit{Point Spread Function}), аппроксимирующей частотно-контрастную характеристику (ЧКХ) линзовой системы:
\begin{equation}
I_{\text{opt}}(x, y) = I_{\text{orig}}(x, y) * H_{\text{PSF}}(x, y)
\end{equation}
где $H_{\text{PSF}}(x, y)$ --- пространственный аналог ЧКХ, моделируемый двумерной функцией Гаусса с дисперсией $\sigma_{\text{opt}}^2$:
\begin{equation}
H_{\text{PSF}}(x, y) = \frac{1}{2\pi\sigma_{\text{opt}}^2} \exp\left( -\frac{x^2 + y^2}{2\sigma_{\text{opt}}^2} \right)
\end{equation}

Матрица фотоприёмника представляет собой дискретную двумерную решётку дельта-функций Дирака $\delta(x,y)$ с пространственными шагами дискретизации $T_x$ и $T_y$ по осям абсцисс и ординат соответственно. Наложение цветоделительного фильтра Байера эквивалентно пространственной модуляции непрерывного сигнала $I_{\text{opt}}(x, y)$ селективными масками каналов $M_c(x, y)$, где $c \in \{R, G_1, G_2, B\}$.

Полная математическая модель дискретизированного мозаичного RAW-сигнала в каждом элементарном отсчёте $I_{\text{RAW}}(x, y)$ выражается как:
\begin{equation}
I_{\text{RAW}}(x, y) = \sum_{n=-\infty}^{\infty} \sum_{m=-\infty}^{\infty} \left[ \sum_{c} I_{\text{opt}}^c(x, y) \cdot M_c(x, y) \right] \cdot \delta(x - nT_x, y - mT_y)
\end{equation}

Для структуры RGGB сенсора Nikon D600 аналитические выражения пространственно-модулирующих масок для каждого канала записываются через тригонометрические базисы:
\begin{equation}
M_G(x, y) = M_{G1}(x,y) + M_{G2}(x,y) = \frac{1}{2} \left[ 1 + \cos\left(\frac{\pi x}{T_x}\right)\cos\left(\frac{\pi y}{T_y}\right) \right]
\end{equation}
\begin{equation}
M_R(x, y) = \frac{1}{4} \left[ 1 + \cos\left(\frac{\pi x}{T_x}\right) \right] \left[ 1 + \cos\left(\frac{\pi y}{T_y}\right) \right]
\end{equation}
\begin{equation}
M_B(x, y) = \frac{1}{4} \left[ 1 - \cos\left(\frac{\pi x}{T_x}\right) \right] \left[ 1 - \cos\left(\frac{\pi y}{T_y}\right) \right]
\end{equation}

\subsection*{2. Спектральный анализ потерь информации и эффект алиасинга}

Переход в пространственно-частотную область с помощью двумерного непрерывного преобразования Фурье ($\mathcal{F}$) позволяет деструктурировать механизм безвозвратной потери спектральных компонентов при классической интерполяции. Умножение на модулирующие гармонические маски $M_c(x,y)$ в координатном пространстве трансформируется в операцию свёртки спектров, что приводит к периодическому сдвигу и копированию исходного непрерывного спектра $\hat{I}(\omega_x, \omega_y)$.

Для зелёного канала (G), имеющего шахматную геометрию, Фурье-образ дискретизированного сигнала принимает вид:
\begin{equation}
\hat{I}_G(\omega_x, \omega_y) = \frac{1}{2}\hat{I}(\omega_x, \omega_y) + \frac{1}{4}\hat{I}\left(\omega_x - \frac{\pi}{T_x}, \omega_y - \frac{\pi}{T_y}\right) + \frac{1}{4}\hat{I}\left(\omega_x + \frac{\pi}{T_x}, \omega_y + \frac{\pi}{T_y}\right)
\end{equation}

Частота Найквиста для решётки зелёных пикселей по диагональным направлениям составляет $\omega_N = \frac{\pi}{\sqrt{2}T}$. Высокочастотные компоненты полезного сигнала, чьи пространственные частоты превышают данный предел, накладываются на боковые спектральные копии (алиасинг), формируя регулярные паразитные интерференционные узоры --- яркостный муар (\textit{zipper effect}).

Для красного (R) и синего (B) каналов шаг дискретизации увеличен в два раза ($2T_x, 2T_y$). Их спектральные функции плотности мощности описываются как:
\begin{equation}
\hat{I}_R(\omega_x, \omega_y) = \frac{1}{4} \sum_{k \in \{0, 1\}} \sum_{l \in \{0, 1\}} \hat{I}\left(\omega_x - \frac{k\pi}{T_x}, \omega_y - \frac{l\pi}{T_y}\right)
\end{equation}

В данном случае спектральные копии смещаются в два раза ближе к центру декартовой частотной системы координат --- на частоты $\left(\frac{\pi}{T_x}, 0\right)$ и $\left(0, \frac{\pi}{T_y}\right)$. В результате области «слепых зон», в которых полезный сигнал полностью перемешан с зеркальными копиями, существенно шире, чем в зелёном канале. Классические линейные дебайеризаторы (билинейная, градиентная интерполяции) восстанавливают эти зоны с помощью низкочастотных фильтров (ФНЧ), что неизбежно приводит к двум деградациям:
\begin{enumerate}
    \item \textbf{Пространственное размытие} (утрата истинной ЧКХ объектива из-за срезания высоких частот);
    \item \textbf{Хроматический муар} (появление ложных низкочастотных цветовых пятен в зонах высокой пространственной активности --- кирпичные кладки, черепица, мелкоструктурные текстуры).
\end{enumerate}

\subsection*{3. Научная новизна и теоретическое обоснование метода}

Научная новизна диссертационного исследования заключается в декомпозиции задачи демозаики как нелинейной обратной задачи восстановления утерянных частотных диапазонов на основе выученного априорного распределения естественных изображений. 

В отличие от известных подходов, ориентированных на минимизацию среднеквадратичного отклонения в пространственной области (L1/L2 метрики), в работе разработан и интегрирован \textbf{гибридный пространственно-частотный критерий оптимизации}. Целевая функция обучения глубокой нейросетевой архитектуры NAFNet модифицирована введением частотного функционала потерь (\textit{Focal Frequency Loss}) на базе двумерного дискретного преобразования Фурье:
\begin{equation}
\mathcal{L}_{\text{total}} = \mathcal{L}_{L1}(I_{\text{pred}}, I_{\text{GT}}) + \gamma \cdot \mathcal{L}_{\text{FFT}}(I_{\text{pred}}, I_{\text{GT}})
\end{equation}
где $\gamma$ --- весовой коэффициент регуляризации частотного домена, а $\mathcal{L}_{\text{FFT}}$ выражается через амплитудные спектры Фурье-образов:
\begin{equation}
\mathcal{L}_{\text{FFT}} = \frac{1}{B \cdot C \cdot H \cdot W} \sum_{b,c,u,v} w(u,v) \cdot \left| \left| \mathcal{F}(I_{\text{pred}})_{b,c,u,v} \right| - \left| \mathcal{F}(I_{\text{GT}})_{b,c,u,v} \right| \right|^2
\end{equation}

Динамическая матрица спектрального взвешивания $w(u,v)$ (\textit{focal weights}) определяется как функция от текущей ошибки аппроксимации:
\begin{equation}
w(u,v) = \left( \frac{\Delta A(u,v)}{\max(\Delta A)} \right)^\alpha, \quad \Delta A(u,v) = \left| \left| \mathcal{F}(I_{\text{pred}})_{u,v} \right| - \left| \mathcal{F}(I_{\text{GT}})_{u,v} \right| \right|^2
\end{equation}

\textbf{Теоретические преимущества разработанного подхода:}
\begin{itemize}
    \item \textbf{Нелинейная частотная интерполяция:} За счет структуры блоков без нелинейных активаций (посредством межканального поэлементного умножения слоев --- \textit{Gated Linear Units}) NAFNet функционирует как адаптивный нелинейный оператор. Модель восстанавливает утерянные гармоники Найквиста не за счет усреднения соседних пикселей, а путем проецирования выученных контекстуальных пространственных связей.
    \item \textbf{Спектральная фокусировка:} Введение динамических весов $w(u,v)$ ($\alpha = 1.0$) заставляет обратный шаг распределения ошибки (вектор градиентов графа вычислений) фокусироваться именно на тех частотных координатах спектра, где алиасинг маски Байера нанес максимальный урон (диагонали Найквиста для G и осевые зоны для R/B).
    \item \textbf{Робастность к физическим шумам тракта:} Моделирование аддитивного белого гауссовского шума, скоррелированного с ядром ЧКХ объектива на этапе синтеза обучающей выборки, позволяет нейросети сформировать устойчивое разделение спектра случайной помехи от спектра истинных высокочастотных контуров изображения.
\end{itemize}

\end{document}
